% Generated by scripts/assemble.py; corrections belong in source/corrections.json.
\examhead{158}{暨南大学 2024 年数学分析试卷}{151}
% p152-b01; source PDF page 152
\par\Needspace{5\baselineskip}
\blockanchor{p152-b01}
\majorquestion{一}计算题.\par
% p152-b02; source PDF page 152
\blockanchor{p152-b02}
\question{1}求极限  $\lim_{n\to\infty}\cos\frac{x}{2}\cos\frac{x}{4}\cdots\cos\frac{x}{2^{n}}(x\neq0)$ .\par
% p152-b03; source PDF page 152
\blockanchor{p152-b03}
\question{2}求极限 $\lim_{x\to 0}\frac{e^{\sin x} - \ln(1 + x) - 1}{[\ln(1 + 2x)]^2}.$\par
% p152-b04; source PDF page 152
\blockanchor{p152-b04}
\question{3}求积分 $\int_0^{2024\pi}\sqrt{1 - \cos 2x} \mathrm{d}x.$\par
% p152-b05; source PDF page 152
\blockanchor{p152-b05}
\question{4}设 $y = x^{2}e^{2x}$，求 $y^{(20)}$.\par
% p152-b06; source PDF page 152
\blockanchor{p152-b06}
\question{5}设 $f(x) = (\arctan x)^{\ln x}$，求 $f'(x)$.\par
% p152-b07; source PDF page 152
\blockanchor{p152-b07}
\question{6}求不定积分 $\int \frac{5x^2 + 1}{x(x^2 + 1)}\mathrm{d}x.$\par
% p152-b08; source PDF page 152
\blockanchor{p152-b08}
\question{7}求定积分 $\int_0^1 x^m (\ln x)^n\mathrm{d}x,m,n\in \mathbb{N}^+$\par
% p152-b09; source PDF page 152
\blockanchor{p152-b09}
\question{8}求极限  $\lim_{n\to\infty}\left(\frac{n}{n^{2}+1^{2}}+\frac{n}{n^{2}+2^{2}}+\cdots+\frac{n}{n^{2}+n^{2}}\right)$ .\par
% p152-b10; source PDF page 152
\blockanchor{p152-b10}
\question{9}求 $2y = x$ 与 $y = x^2$ 围成区域绕 $y$ 轴旋转所得旋转体的体积.\par
% p152-b11; source PDF page 152
\par\Needspace{5\baselineskip}
\blockanchor{p152-b11}
\majorquestion{二}证明题.\par
% p152-b12; source PDF page 152
\blockanchor{p152-b12}
\question{1}已知 $a_{n}, b_{n}$ 是正数列, $k_{n} = \frac{b_{n} a_{n}}{a_{n+1}} - b_{n+1}$, 存在 $c > 0$, 使得对任意的正整数 $n$, 均有 $k_{n} > c$, 证明: 级数 $\sum_{n=1}^{\infty} a_{n}$ 收敛.\par
% p152-b13; source PDF page 152
\blockanchor{p152-b13}
\question{2}设 $f_0, f_1$ 可微, 且 $f_0 \neq 0$, $\left| \begin{array}{cc} f_0 & f_1 \\ f_0' & f_1' \end{array} \right| \neq 0$, 证明: 对任意的 $\lambda_0, \lambda \in \mathbb{R} \setminus \{0\}$, 函数 $F(x) = \lambda_0 f_0(x) + \lambda_1 f_1(x)$ 至多存在一个零点, 且若存在零点, 则该点 $F(x)$ 的一阶导数不为零.\par
% p152-b14; source PDF page 152
\par\Needspace{6\baselineskip}
\blockanchor{p152-b14}
\question{3}设  $P(x,y)$ ,  $Q(x,y)$  可微, 区域 D 是由曲线 L:  $\left\{\begin{aligned} x &= \varphi(t) \\ y &= \psi(t) \end{aligned}\right.$  所围成的区域, 且  $P_{x} + Q_{y} = 0$  无根.
证明:\par
\nopagebreak[4]
% p152-b15; source PDF page 152
\blockanchor{p152-b15}
\displayformula{\left\{ \begin{array}{l} \varphi^ {\prime} (t) = P (\varphi (t), \psi (t)); \\ \psi^ {\prime} (t) = Q (\varphi (t), \psi (t)). \end{array} \right.}
% p152-b16; source PDF page 152
\blockanchor{p152-b16}
一定不成立.\par
% p152-b17; source PDF page 152
\blockanchor{p152-b17}
\question{4}设 $f(x)$ 在 $[a, b]$ 上连续，且 $\int_{a}^{b} f^{2}(x) \mathrm{d}x = 0$，证明：$f(x)$ 在 $[a, b]$ 上恒为零。\par
% p152-b18; source PDF page 152
\par\Needspace{5\baselineskip}
\blockanchor{p152-b18}
\majorquestion{三}分析题.\par
% p152-b19; source PDF page 152
\blockanchor{p152-b19}
\question{1}分析 $f(x) = x^{2} \cos \left( \frac{\pi}{2} - x \right) - x^{3} + x^{4}$ 在 0 处是否取得极值.\par
% p152-b20; source PDF page 152
\blockanchor{p152-b20}
\question{2}设 $f_{n}(x) = e^{n(x - 1)}, x \in (0,1)$, 分析 $\{f_{n}(x)\}$ 在 $(0,1)$ 上的收敛性和一致收敛性.\par
